\documentclass[10pt]{article}
\usepackage[margin=20mm]{geometry}
\usepackage{amsmath,amssymb,mathtools,fontspec,unicode-math,xurl,hyperref}
\IfFontExistsTF{DejaVu Serif}{\setmainfont{DejaVu Serif}}{\setmainfont{Latin Modern Roman}}
\setmathfont{latinmodern-math.otf}
\setlength{\parindent}{0pt}\setlength{\parskip}{6pt}\sloppy
\title{Technical verification of
Positive curvature, torus symmetry,
and matroids}
\author{AI-assisted technical review}
\date{Focused independent acceptance}
\begin{document}
\maketitle
Abstract: This review concerns the 56-page arXiv:2212.08152v2 source. The frozen initial review read the full text; the present independent acceptance rechecks its three reported Problems. The generalized small-cycle proof contains a false connectivity assertion, disproved by an admissible three-connected cubic graph. The actual later applications inspected use at most three deleted cycle edges, for which connectivity holds; a replacement also proves the full stated estimate. Two further issues concern normalized equality witnesses and fixed versus optimized representation weights. All three are assessed Minor within these precise limits. Initial exact positive calculations are retained with their original attribution. No counterexample to Theorems A–C is established. Scientific adjudication remains pending.\par
\section*{Source identity and coverage}
The paper is by Lee Kennard, Michael Wiemeler and Burkhard Wilking, associated with JEMS DOI 10.4171/jems/1656 (https://doi.org/10.4171/jems/1656). The inspected file is arXiv:2212.08152v2 (https://arxiv.org/abs/2212.08152v2). Its side label says 10 June 2025 and its internal date says 12 June 2025. The actual PDF has 56 pages; the metadata comment’s 55-page description is not used to override that count. These are preserved version details, not evidence of a scientific defect. All locations below use one-based PDF pages, agreeing with printed numbers.\par
The source manifest binds this exact author manuscript. Its equivalence to the journal version of record has not been established; the normal publisher PDF route returned a subscription page. The report therefore makes no claim about an uninspected publisher revision. Acquisition searches and an additional bounded publisher/arXiv search did not locate separate code, supplements or an erratum. That is not an exhaustive absence claim. No author software was executed, and lack of code is not counted against the paper.\par
The frozen initial review read Sections 1–5 and the references. Its coverage is retained as historical; the present acceptance focuses on the reported Problems and their actual dependencies. Detailed checks are identified below. The graph-embedding case analysis on pp. 23–26, the external obstruction lists, Seymour’s theorem, and the imported connectedness, splitting and periodicity theorems were not independently reproved. The geometrical proof was traced through its stated dependencies, not formally certified.\par
\section*{Scope of the principal conclusions}
Theorem A assumes a closed, oriented, positively curved manifold, an effective isometric torus action with a fixed point, and odd numbers of isotropy components near that point. Its conclusions concern rational cohomology, globally for rank nine and through degree \(n/3\) for rank six. It does not assert diffeomorphism classification or remove the isotropy hypothesis.\par
Theorems B and C provide graph systole and regular-matroid cogirth bounds through rank nine. Their tables differ at rank five: \(s(5)=3/8\) whereas \(c(5)=2/5\), because the sporadic matroid contributes. Confusing these two optimizations would lose an essential distinction. The route to Theorem A passes through the codimension bounds, the six-involution refinement, and the periodicity arguments. Section 5.4 supplies a separate shorter route for rank nine; it still uses the matroid bounds and geometric input theorems.\par
\section*{Findings}
Problem 1 (Deletion of cycle edges need not preserve connectivity).  Minor: false general-range proof assertion, with the inspected downstream range unaffected.\par
Lemma 2.3 on p. 14 asserts \[\operatorname{sys}(G)^{-1}\geq \frac hg+s(b-h)^{-1},
 \qquad 1\leq h\leq\min(g,b-1),\] for a three-connected graph of Betti number \(b\) containing a \(g\)-cycle. Its proof handles one deleted edge, then deletes the \(h\) heaviest edges of the cycle. The last paragraph asserts that the graph remains connected, even after deleting the whole cycle, on the basis of three-connectivity. This inference is false.\par
\subsection*{A witness satisfying the actual range}
Take the ordinary cube graph with vertices the three-bit strings, encoded by integers \(0,\ldots,7\), and edges joining strings that differ in one bit. It is cubic, three-edge-connected and three-vertex-connected; the independent finite check verifies both connectivity conventions after every deletion of at most two edges or vertices. It has \(12\) edges, \(8\) vertices and \(b=5\). Choose \[C=(0,1,3,2,6,7,5,4,0),\qquad g=8,\qquad h=4.\] Thus \(h\leq\min(8,5-1)\), exactly as required. Let \[A=\{01,23,45,67\}\subset E(C).\] Give edges in \(A\) weight \(2/16\) and the remaining eight edges weight \(1/16\). The total is one and the four edges in \(A\) are the strictly heaviest edges of \(C\), eliminating any tie-breaking concern.\par
Deleting \(A\) leaves the two four-cycles on the even and odd vertices. The graph is disconnected, and \[b_1(G\setminus A)=8-8+2=2,\qquad b-h=1.\] The precise connectivity and Betti-number premise used to generalize the proof therefore fails on an admissible input. The witness does not dispute the numerical inequality.\par
\subsection*{Replacement derivation and impact}
Let \(t\) be the shortest cycle length for a normalized nonnegative weighting, let \(A\) be the \(h\) heaviest edges of the chosen cycle, and set \(a=\lambda(A)\). Then \[a\geq (h/g)\lambda(C)\geq (h/g)t.\] Even if \(H=G\setminus A\) disconnects, its Betti number is at least \(b-h\geq1\). Delete further non-bridge edges as needed to obtain a subgraph \(H'\) of Betti number exactly \(b-h\). Its total weight \(L\) is at most \(1-a\), and all its cycles remain cycles of \(G\). If \(H'\) is disconnected, join its components by zero-weight bridge edges. This operation preserves the Betti number, total weight and cycles, exactly as in the first step of Lemma 2.2 on p. 13. Thus the connected-graph convention in the systole definition causes no difficulty. The inductive bound at Betti number \(b-h\) gives \[t\leq Ls(b-h)\leq(1-a)s(b-h)
       \leq \bigl(1-(h/g)t\bigr)s(b-h).\] If \(L=0\), a zero-weight cycle already gives the desired conclusion; otherwise the normalization is legitimate. Rearrangement proves the stated estimate without the false connectivity claim.\par
The affected range must be separated from the actual uses. Page 16 applies the estimate with \(h=1\) and with \(h=2\) when \(b=10,g=3\); strict monotonicity also follows from \(h=1\). Page 21 uses \(h=g=3\) for \(b=9\). The two cases of Lemma 2.8 on p. 22 delete pairs of edges from a four-cycle, rather than deleting that whole cycle.\par
For any three-edge-connected graph and any \(A\subseteq E(C)\) with \(|A|\leq3\), disconnection would give a nonempty cut \(\delta(S)\subseteq A\). Three-edge-connectivity forces \(|\delta(S)|=3\), whereas a cycle meets a cut in an even number of edges. This contradiction proves connectivity in precisely the inspected downstream cases. The cube example with \(h=4\) therefore exposes the overextended proof assertion, but does not invalidate those applications or the table and geometric conclusions. The Minor assessment rests on this actual dependency boundary, not on the shortness of a repair. The replacement is still required for the lemma’s full advertised range.\par
Problem 2 (The displayed equality weights are not normalized).  Minor: local equality-witness corrections, with independently verified replacements.\par
On p. 19, the \(F_{12}\) calculation uses orbit sizes \(8,8,2\) and the normalization \[8\lambda+8\mu+2\nu=1.\] Its displayed equality choice is \[(\lambda,\mu,\nu)=(1,2,3)/28.\] This has total weight \(30/28=15/14\), so is not admissible for the optimization as defined on p. 12. The printed three cycle lengths are \(8/28\), \(9/28\) and \(8/28\). After normalizing the total weight, the shortest cycle has length \(4/15\), strictly less than \(2/7\). Thus the stated triple does not certify the asserted equality.\par
Replacing the last numerator by two gives \[(\lambda,\mu,\nu)=(1,2,2)/28.\] It has total weight one and makes the three displayed cycle constraints equal to \(2/7\). More strongly, an independent exhaustive enumeration of all \(104\) simple cycles of the graph reconstructed from Figure 2.2 verifies that every cycle has length at least \(2/7\). Together with the paper’s averaging upper bound, this proves \(\operatorname{sys}(F_{12})=2/7\) with the corrected weights.\par
The reconstruction uses the printed vertices \(1,\ldots,9\); label the unnumbered upper-left vertex \(10\) and the upper/lower inner-right vertices \(11,12\). The four edges in the left square and four in the right square have weight \(\mu\); the eight connecting diagonal edges have weight \(\lambda\); the top and bottom horizontal edges have weight \(\nu\). The exact adjacency and all cycle records are saved with the calculation, so the finite model can be audited against the source figure.\par
Similarly, p. 16 assigns weight \(1/4\) to all six edges of \(K_4\) while working with total weight one. The normalized value is \(1/6\). Enumerating its seven simple cycles gives minimum \(1/2\) with the corrected weighting. Unlike \(F_{12}\), merely rescaling the erroneous uniform \(K_4\) weighting already preserves the desired ratio. These are local witness defects: the \(F_{12}\) upper bound does not use the erroneous equality triple, and the corrected attainments support the stated constants.\par
Problem 3 (A fixed representation is confused with a maximum over weights).  Minor: false displayed equality in the explanatory reduction.\par
On p. 12 the correct identity for fixed multiplicities is followed by \[c(\rho)=\max_\lambda\min_C\lambda(C).\] The left side depends on the given representation; the right side is the optimized invariant of the underlying graph. The equality does not hold for arbitrary multiplicities.\par
For a concrete admissible example, take the cographic theta representation of \(T^2\), with weights \((1,0),(0,1),(1,1)\) and multiplicities \(1,1,10\). Its isotropy groups are connected: each individual weight is primitive and every independent pair has determinant \(\pm1\); repetitions do not change the generated lattice. The circle in direction \((1,-1)\) fixes the ten copies of the third weight, so \[c(\rho)=\frac{2}{12}=\frac16,
       \qquad \operatorname{sys}(\Theta)=\frac23.\] The three possible annihilator lines give moved-weight ratios \(11/12,11/12,1/6\); a generic direction moves all weights. This verifies the minimum. Equal theta edge weights attain \(2/3\), and averaging its three two-edge cycle sums gives the matching upper bound.\par
For a fixed representation, write \(c(\rho)\leq\operatorname{sys}(G)\), or put the maximum over multiplicities on the left as well. Section 3.2 correctly distinguishes \(c(H,\lambda)\) from \(c(H)\), which is substantive counterevidence against interpreting the earlier display as a mistaken global strategy. The upper-bound reduction requires only the corrected inequality.\par
\section*{Positive checks retained from the initial review and remaining limits}
The initial review used the source figures and formulas on pp. 19–21 to construct finite graphs independently, without executing author software. Its preserved exact rational arithmetic gave:\par
Object | Edges | Simple cycles | Verified normalized minimum\par
\(K_4\), corrected weights | 6 | 7 | \(1/2\)\par
\(F_{12}\), corrected weights | 18 | 104 | \(2/7\)\par
\(F_{13}\), printed weights | 18 | 109 | \(8/27\)\par
\(F_{14}\), printed weights | 18 | 109 | \(3/10\)\par
The \(F_{13}\) model is its nine-cycle with the three repeated tripod labels identified as in Figure 2.3. The \(F_{14}\) model includes both extra edges specified in the caption of Figure 2.4, although not drawn there. Omitting those caption edges would test the wrong graph.\par
For \(R_{10}\), all \(31\) nonzero binary dual vectors were evaluated on the displayed ten columns. Minimum support size was four, so uniform weights give \(2/5\). The five adjacent-coordinate dual vectors used in Proposition 3.11 each have four nonzero evaluations and cover each column twice, verifying the averaging certificate as well as the lower witness. These checks support the inspected local constants; they do not verify Seymour’s decomposition theorem or the entire rank-nine graph classification.\par
The geometric reading checked the distinction between circle and involution fixed-point sets, the use of the odd-isotropy hypothesis to obtain circle fixed components, parity of codimensions, and the rank-nine inequalities in Section 5.4. Imported hypotheses were tracked, but no independent proof of the connectedness or four-periodicity theorems is claimed. Minor reciprocal/index slips elsewhere were not expanded into additional findings without a distinct effect.\par
The technical corrections are to replace the general-range proof step in Problem 1 and correct the explicit local statements in Problems 2–3. The present independent acceptance supports these three Minor findings within the boundaries stated above. The verified counterexamples concern those statements, not the final classification; scientific adjudication remains pending.\par
\section*{Supporting checks}
Paper: `M25-0283`. Its original report, complete-reading ledger, source-image views, independent script, exact graph inputs, cycle records, execution receipts and build history remain preserved. They are historical work, not a new complete-paper reading.\par
The assigned acceptance assistant independently checked the three Problems, their admissible witnesses, counterevidence, repair arguments, source locations and downstream use. It compared the \(F_{12}\) model directly with Figure 2.2, reviewed the exact calculation and its retained execution binding, and made the present severity and scope decisions. Problem 1 is revised from Major to Minor for the specific actual-use boundary above. The native bounded read-only assistant `/root/check\_0280\_0281` cross-checked the graph transcription, cycle records, edge orbits and small-cycle applications.\par
The acceptance assistant revised this report and performed the actual offline compilation, PDF parsing and inspection of every current report page. Current source, TeX, PDF, compiler and dependency hashes, retained payloads, logs and page-view records are bound in the provenance and delivery inventory. Original false render records remain unchanged; rendering and actual viewing are separate facts.\par
The controller’s UI observations retain their controller attribution. Runtime instructions identify GPT-6 Astra Pro; backend correspondence was not independently verified. No author software was run. Exact author-version limits and external proof limits remain in force. Scientific adjudication is `pending`; `human\_reviews=[]`.\par
\end{document}
